\chapter{系统组成及功能说明}
阐述具体的设计细节，采用图文结合的方式。
\section{整体介绍}
给出系统整体框图，清楚标注各子模块，并用文字说明模块之间的关系、数据流、控制流和关键接口。
\begin{figure}[htbp]
\centering
\fbox{\parbox[c][32mm][c]{.86\textwidth}{\centering 系统整体框图占位\\请替换为作品的实际框图}}
\caption{系统整体框图（待替换）}
\label{fig:system}
\end{figure}
\section{硬件系统介绍}
\subsection{硬件整体介绍}
说明 FPGA 器件、开发平台、时钟、复位、电源、外设和接口，以及硬件各部分的关系。
\subsection{机械设计介绍}
如果有机械结构，从总体到局部逐级给出各组件的具体设计图，可采用 CAD 文件截图或者手绘图片；若无机械结构，明确说明不涉及。
\subsection{电路各模块介绍}
从总体到局部逐级给出各模块的具体设计图，并标记关键输入、输出信号线。可采用电路图、SCH 原理图、PCB 版图等截图。说明 FPGA 各逻辑模块的接口、时序和连接关系。
\section{软件系统介绍}
\subsection{软件整体介绍}
结合关键图片介绍软件体系，包括 PC 端或云端（如有）、FPGA 逻辑、固件以及软硬件交互方式。
\subsection{软件各模块介绍}
根据总体框图给出各模块的具体设计说明。从顶层到底层逐次给出各函数的流程图及关键输入、输出变量；对于 FPGA 逻辑，可结合状态机、数据通路、时序图和模块接口说明。
